% TOP_PROBLEM: {"id": 30003025, "title": "Color Savings and Strong Chromatic Index Bounds", "category_id": 3, "category": {"id": 3, "name": "graph_theory", "display_name": "Graph Theory", "description": "Problems involving graphs, networks, and their properties.", "slug": "graph-theory", "order_index": 3, "created_at": "2026-07-31T15:26:25.671Z"}, "status": "open"}
% FIRST_POSTED: 2026-10-06
% ATTEMPT_STATUS: unresolved
% OPENAI_CLAIM: {"schema_version":1,"source_repo":"https://github.com/openai/math","source_commit":"adc7f1241b42e322a6451854ab7e4b4c146bf78a","release_date":"2026-10-06","match":"partial","resolution":"partial","independently_reviewed":false,"second_pass":"n/a","families":[{"family":"184","title":"Correspondence coloring with a fixed forbidden subgraph","subject":"Combinatorics","manuscripts":[{"title":"Correspondence coloring graphs with a forbidden clique","pdf_path":"preprints/Correspondence-Coloring-Graphs-with-a-Forbidden-Clique-October-5-2026/correspondence-coloring-forbidden-clique.pdf","readme_path":"preprints/Correspondence-Coloring-Graphs-with-a-Forbidden-Clique-October-5-2026/README.md","date":"2026-10-05","theorem_ref":"Theorem 1.1"}],"lean":{"doc_path":"lean/docs/184.md","comparators":[],"declarations":[],"covers_main_theorem":false},"reasoning_trace":null}]}
\documentclass[11pt]{article}
\usepackage[margin=25mm]{geometry}
\usepackage{fontspec}
\setmainfont{Latin Modern Roman}
\usepackage{amsmath,amssymb,mathtools}
\usepackage{unicode-math}
\setmathfont{Latin Modern Math}
\usepackage{enumitem}
\usepackage{xurl,hyperref,bookmark}
\hypersetup{colorlinks=true,urlcolor=blue}
\setcounter{secnumdepth}{0}
\setlength{\parindent}{0pt}
\setlength{\parskip}{5pt}
\setlength{\emergencystretch}{3em}
\begin{document}
\section*{Color Savings and Strong Chromatic Index Bounds}

\subsection{OpenAI's claim}
\textit{Correspondence coloring with a forbidden clique}: For each fixed $r\ge4$, the paper claims correspondence chromatic number $O_r(\Delta/\log\Delta)$ in graphs with no $K_r$ subgraph and large maximum degree. Ordinary and list-coloring versions follow, and fixed forbidden subgraphs are also covered.

\subsection{Outline of the argument}
\textit{Correspondence coloring with a forbidden clique}: Random partial-coloring rounds act directly on the graph of available list entries and forbidden matching pairs. Weights track total list supply, individual probabilities, and conflicting edge mass. An entropy-and-mass potential pays for losses when neighboring vertices are colored, including local adjustments for triangles, until the residual lists can be completed.

\subsection{Scope relative to this problem}
\textit{Correspondence coloring with a forbidden clique}: The theorem gives a quantitative color saving for ordinary, list and correspondence coloring when the clique number is bounded and maximum degree is large; the general degree/clique tradeoff and strong edge-coloring clause remain unsettled.

This page records partial support. Complete resolution of the catalogue question is not established by this record.

\subsection{Formal verification}
Family 184 describes Lean coverage for a large independent set in a fixed-clique-free graph at the logarithmically improved scale; correspondence coloring is separate.

\textit{Correspondence coloring with a forbidden clique}: No manuscript-specific Comparator is listed in the pinned inventory.

The documented selection does not cover the entire manuscript assertion relevant here.

This site has not run the Lean development. These links describe OpenAI’s selected formal statements and should be read with their stated scope.

\subsection{Status caveat}
OpenAI reports that an internal model produced the manuscripts. Their arguments have not been independently verified as proofs by this site. Any portion outside the documented formal scope may contain mistakes, and the catalogue label reports the announced mathematical claim.

\subsection{Sources}
\begin{itemize}
\item OpenAI, \emph{Correspondence coloring graphs with a forbidden clique}, 2026-10-05; citation key \texttt{OAI:Correspondence-Coloring-Graphs-with-a-Forbidden-Clique-October-5-2026}. \href{https://github.com/openai/math/blob/adc7f1241b42e322a6451854ab7e4b4c146bf78a/preprints/Correspondence-Coloring-Graphs-with-a-Forbidden-Clique-October-5-2026/correspondence-coloring-forbidden-clique.pdf}{PDF (GitHub)}; \href{https://github.com/openai/math/raw/adc7f1241b42e322a6451854ab7e4b4c146bf78a/preprints/Correspondence-Coloring-Graphs-with-a-Forbidden-Clique-October-5-2026/correspondence-coloring-forbidden-clique.pdf}{PDF (direct)}; \href{https://github.com/openai/math/blob/adc7f1241b42e322a6451854ab7e4b4c146bf78a/preprints/Correspondence-Coloring-Graphs-with-a-Forbidden-Clique-October-5-2026/README.md}{Manuscript page and citation}; \href{https://github.com/openai/math/blob/main/preprints/Correspondence-Coloring-Graphs-with-a-Forbidden-Clique-October-5-2026/correspondence-coloring-forbidden-clique.pdf}{Latest version}.
\item Family 184: \href{https://github.com/openai/math/blob/adc7f1241b42e322a6451854ab7e4b4c146bf78a/lean/docs/184.md}{Lean scope}.
\item \href{https://github.com/openai/math/blob/adc7f1241b42e322a6451854ab7e4b4c146bf78a/CONTENTS.md}{OpenAI catalogue entry}.
\end{itemize}
\end{document}
